% ---------------------------------------------------------------------
% STATUT : RÉDIGÉ. Chapitre central de la version actuelle du mémoire.
% Tous les chiffres proviennent de scripts vérifiés :
%   6_contagion/30 (bornes), 31 (comparaison des régimes), 33 (relatif/absolu).
% Figures : Z_identification_partielle, Z2_comparaison_regimes, Z4_relatif_vs_absolu.
% ---------------------------------------------------------------------
\chapter{Identification partielle et bornes de capital}
\label{chap:partielle}

\begin{cle}
\textbf{Le problème, et la sortie.} Le chapitre précédent établit que la \emph{direction}
de la contagion n'est pas identifiable sur les données disponibles. Deux attitudes sont
alors possibles. On peut \textbf{poser} $W$ au jugement d'expert et publier un SCR
ponctuel : c'est simple, lisible, et entièrement suspendu à un jugement qui n'a pas encore
été élicité. Ou bien on peut refuser de choisir, caractériser l'\textbf{ensemble} des $W$
compatibles avec ce que l'on sait, et publier des \textbf{bornes}. Ce chapitre construit
la seconde voie, la compare à la première, et montre qu'elle est la seule des deux à être
opposable aujourd'hui.
\end{cle}

\section{Le modèle en couches}

Le modèle se lit désormais en deux couches, dont la séparation est la ligne de partage
entre ce qui est estimé et ce qui est posé.

\paragraph{Le socle, identifié.} Fréquence d'entrée surdispersée, accumulation par cause
commune, sévérité GPD en euros. Aucun de ces éléments ne dépend de $W$. En annulant toute
contagion ($W = 0$), on obtient le \textbf{SCR nu} : $5\,275$~M\euro{} dans le calage de
référence. C'est le plancher défendable, et il ne repose sur aucun paramètre non
identifiable.

\paragraph{La couche de contagion, partiellement identifiée.} Toute matrice se décompose
de façon unique en
\[
  W \;=\; S + A, \qquad S = \tfrac{1}{2}(W + W^{\top}), \qquad A = \tfrac{1}{2}(W - W^{\top}) .
\]
La donnée identifie la co-occurrence, donc $S$. Elle n'identifie pas la direction, donc
$A$ (chapitre~\ref{ch:identifiabilite}). L'\textbf{ensemble admissible} est alors
\[
  \mathcal{W}(t) \;=\; \bigl\{\, S + A \;:\; A^{\top} = -A, \;\; |A_{jk}| \leq t\,S_{jk},
  \;\; \rho(S+A) < 1 \,\bigr\},
\]
la contrainte $|A_{jk}| \leq S_{jk}$ découlant de la seule positivité ($W_{jk}\geq 0$ et
$W_{kj}\geq 0$), et $t\in[0,1]$ mesurant ce que la donnée \emph{ignore} de la direction :
$t=1$ est l'ignorance totale, c'est-à-dire l'état actuel des sources publiques.

\section{Ce que le capital devient}

Les bornes sont obtenues par énumération \textbf{exhaustive} des $2^{10} = 1\,024$ sommets
du pavé, ce qui les rend déterministes et reproductibles. À $t=1$ :
\[
  \mathrm{SCR} \;\in\; [\,6\,858 \;;\; 8\,697\,]~\text{M\euro{}} ,
\]
soit un facteur $1{,}27$ entre les deux bouts. Rapporté au socle, cela s'énonce ainsi :
\textbf{la contagion ajoute entre $+30\,\%$ et $+65\,\%$ au capital, et la donnée ne
permet pas de trancher plus finement}. Le point d'expert ($8\,110$~M\euro{}) est un point
de cet ensemble, jamais sa mesure.

\begin{attention}
\textbf{Une lecture à ne pas faire.} À indice de queue élevé, la queue annuelle obéit au
principe de la perte unique dominante : la VaR est \emph{peu} sensible à la structure de
cascade. Faire porter la démonstration par le seul niveau de VaR donnerait un tableau de
sensibilité plat et trompeusement rassurant. La contribution se lit sur la \textbf{perte
moyenne} (bornes $[849\,;1\,100]$~M\euro{}, estimateur bien moins bruité) et surtout sur
la \textbf{décision}, c'est-à-dire l'ordre de remédiation.
\end{attention}

\paragraph{La valeur de l'information.} Le paramètre $t$ se lit comme le degré de
contrainte qu'un registre d'incidents apporterait sur la direction. Passer de $t=1$ à
$t=0{,}5$, c'est-à-dire borner l'asymétrie à la moitié du niveau de co-occurrence, ferait
tomber la largeur des bornes de $1\,839$ à $695$~M\euro{}, soit $\mathbf{-62\,\%}$. La
spécification de reporting proposée au chapitre~\ref{ch:identifiabilite} cesse ainsi
d'être un v\oe{}u : elle est chiffrée en euros de capital.

\section{Poser \texorpdfstring{$W$}{W} ou le borner : la comparaison avant élicitation}

Le test décisif se fait \emph{avant} les séances d'élicitation, car un dispositif qui ne
survit qu'à une élicitation réussie n'est pas un dispositif, c'est un pari.

\begin{center}\footnotesize
\renewcommand{\arraystretch}{1.3}
\begin{tabular}{@{}p{4.1cm} p{4.5cm} p{5.6cm}@{}}
\toprule
 & \textbf{Régime A : on pose $W$} & \textbf{Régime B : modèle en couches} \\
\midrule
Sortie sur le capital & un point, $8\,110$~M\euro{} & socle $5\,275$ + bande $[6\,858\,;8\,697]$ \\
Opposable sans élicitation & non, tout repose sur le classeur & oui, $79\,\%$ du capital est déjà fixé \\
Panel biaisé & déplace le SCR de $777$~M\euro{} & la bande ne bouge pas \\
Priorité de remédiation & P1, sans marge d'erreur & tête P1/P4, ordre interne ouvert \\
Si l'élicitation échoue & plus de paramètre, plus de modèle & le résultat tient, en bornes \\
\bottomrule
\end{tabular}
\end{center}

\paragraph{Le test de biais du panel.} On fait varier la direction élicitée depuis
\og{}le panel confirme le classeur\fg{} jusqu'à \og{}le panel dit l'inverse\fg{}, en
allant au bord du pavé admissible. Le régime~A suit : le SCR publié se déplace de
$777$~M\euro{}, soit $10\,\%$, et la priorité bascule entre trois piliers différents,
\textbf{sans que rien dans le régime~A ne signale ce déplacement au lecteur}. La bande du
régime~B, elle, ne bouge pas d'un euro, et contient tous les dérapages possibles du panel.

\paragraph{À quoi sert alors l'élicitation.} Son rôle est précis et limité, et c'est un
résultat acquis sans elle. Sur l'ensemble admissible, \textbf{P5 n'arrive jamais en
tête}, P3 dans moins de $10\,\%$ des cas, et P1 ($45\,\%$) et P4 ($36\,\%$) concentrent
$81\,\%$ des configurations. L'élicitation sert donc à \textbf{départager P1 et P4} et à
resserrer la bande. Elle ne peut pas faire remonter les autres piliers : l'ensemble
admissible l'interdit déjà.

\section{Le niveau est illustratif, le rapport est le résultat}

Un relecteur exigeant objectera qu'une tour d'hypothèses accouche d'un grand nombre.
L'objection est fondée si le mémoire vend un niveau. Elle tombe si le mémoire vend un
rapport \emph{et} démontre que le rapport, lui, ne bouge pas. On fait donc varier ce que
le modèle ne maîtrise pas, en laissant la cascade inchangée.

\begin{center}\footnotesize
\renewcommand{\arraystretch}{1.25}
\begin{tabular}{@{}p{4.3cm} r r r r@{}}
\toprule
\textbf{Configuration} & \textbf{socle} & \textbf{bornes} &
\textbf{haute/socle} & \textbf{\% fixé} \\
\midrule
Référence (OpRisk)   & $5\,275$  & $[6\,858\,;8\,697]$   & $1{,}65$ & $78{,}9\,\%$ \\
Sévérité PRC         & $239$     & $[336\,;424]$         & $1{,}77$ & $79{,}2\,\%$ \\
Fréquence $\times 3$ & $11\,709$ & $[14\,353\,;17\,405]$ & $1{,}49$ & $82{,}5\,\%$ \\
\bottomrule
\end{tabular}
\end{center}

\paragraph{Lecture.} Le \textbf{niveau} bouge d'un facteur $41$. Il est fixé par la
calibration de sévérité et de fréquence, qui sont des \emph{entrées} du modèle et non des
sorties de la cascade, et qui dépendent d'un périmètre de collecte non transposable
(pertes américaines de grandes institutions, plafond de réassurance). Les
\textbf{rapports} varient d'au plus $33\,\%$ de leur valeur moyenne, soit un ordre de
grandeur trente fois plus faible, et le \textbf{groupe de tête} de la priorité de
remédiation est identique dans les trois configurations.

\begin{attention}
\textbf{Où la stabilité du relatif s'arrête, et il faut le dire.} Le rapport le moins
stable est la largeur de bande rapportée au socle, qui varie de $32{,}9\,\%$ : elle
diminue quand la fréquence augmente, ce qui est attendu, la queue annuelle étant alors
moins dominée par une perte unique et donc moins sensible à la structure de cascade. La
stabilité du relatif n'est donc pas parfaite. Elle est simplement d'un tout autre ordre
que celle du niveau.
\end{attention}

\begin{cle}
\textbf{Ce que ce chapitre autorise à écrire.} Le niveau absolu du SCR est
\textbf{illustratif} et ne doit pas être lu comme une mesure. Ce que le modèle détermine,
et qui résiste au changement de source, c'est le rapport du capital au socle, la largeur
de la bande d'ignorance directionnelle, et le classement des piliers. Ce n'est pas une
précaution de langage ajoutée après coup : c'est un résultat vérifié, et c'est la raison
pour laquelle le mémoire présente partout des écarts et des bornes plutôt que des points.
\end{cle}

\begin{figure}[htbp]
\centering
\includegraphics[width=\linewidth]{Z_identification_partielle.png}
\caption{Identification partielle de la contagion dirigée. (a)~Le coût de la donnée
manquante : les bornes se resserrent à mesure que l'ignorance directionnelle $t$ diminue.
(b)~La lecture centrale du prior d'expert, à l'intérieur de bornes qui n'en dépendent pas.
(c)~La priorité de remédiation robuste, au sens du regret maximal. (d)~La dominance
majoritaire entre piliers sur l'ensemble admissible.}
\label{fig:partielle}
\end{figure}

\begin{figure}[htbp]
\centering
\includegraphics[width=\linewidth]{Z2_comparaison_regimes.png}
\caption{Poser $W$ ou le borner, comparés avant élicitation. (a)~Un point contre un socle
et des bornes. (b)~Seule la part supérieure dépend de l'élicitation. (c)~Le régime~A suit
le biais du panel, la bande du régime~B ne bouge pas. (d)~Ce que l'élicitation doit
départager, et rien de plus.}
\label{fig:regimes}
\end{figure}

\section{L'échelle de repli sur la direction}
\label{sec:echelle-u}

La décomposition $W = S + A$ borne la direction en bloc. On peut faire mieux, en suivant
la méthode d'attaque puis de repli : reparamétrer la propagation par
\[
  \operatorname{logit} p_{ij} \;=\; \operatorname{logit} p_j + u_{ij},
\]
où $p_j$ est l'\textbf{attractivité de la cible} (ce qui subsiste quand la source ne compte
pas) et $u_{ij}$ l'\textbf{excès propre à la source}, c'est-à-dire la direction. La
décomposition n'est identifiée qu'une fois posée la normalisation $\sum_{i\neq j} u_{ij}=0$
pour chaque $j$ : c'est une convention de repérage, et non une information, ce que l'on
vérifie numériquement (une translation compensée de $p_j$ et de $u$ laisse le SCR
strictement inchangé). On descend alors une échelle de crans, du plus ambitieux au repli :
$u$ libre, puis $u$ de \textbf{rang $1$} ($u_{ij}=b_j(a_i-\bar a)$, qui capte $74\,\%$ de
la direction posée et resserre les bornes de $25\,\%$ au prix d'une hypothèse de structure),
puis $u$ élicité, puis $u=0$.

\begin{cle}
\textbf{Le cran nul, et ce qu'il révèle de la contribution.} À $u=0$, la propagation ne
dépend plus que de la cible : \emph{aucune source ne se distingue}. Or c'est exactement là
que P1 et P4, les deux piliers de tête, deviennent \textbf{indiscernables} (gains de
remédiation $1\,051$ contre $1\,052$~M\euro{}, écart nul), alors qu'ils sont séparés de
$353$~M\euro{} dès que la direction du classeur est rétablie. \textbf{La direction est donc
précisément ce qui départage les deux piliers prioritaires} : sans elle, on ne sait plus
lequel remédier d'abord. La contribution du mémoire ne se lit pas sur le niveau d'un
quantile, elle se lit ici, sur une décision qui bascule (figure~\ref{fig:echelle}).
\end{cle}

\begin{figure}[htbp]
\centering
\includegraphics[width=\linewidth]{Z5_echelle_repli.png}
\caption{L'échelle de repli sur $u_{ij}$. (a)~Les bornes de capital se resserrent quand on
recule d'un cran (hypothèse de structure). (b)~Un rang $1$ capte $74\,\%$ de la direction
posée. (c)~Ce que la direction change à la décision : au cran nul, les gains de remédiation
de P1 et P4 se rejoignent.}
\label{fig:echelle}
\end{figure}

% =====================================================================
